\usepackage[scr=rsfs,cal=euler]{mathalfa}
\hyphenation{admit-ting embed-ded immer-sion using}

\newcommand{\Changel}{\mathcode`l="7160}
\newcommand{\Changelback}{\mathcode`l="716C}
\newcommand{\NoChangel}[1]{%
\expandafter\let\csname old\string#1\endcsname=#1
\let#1=\relax
\newcommand{#1}{\mathop{\mathcode`l="716C\csname old\string#1\endcsname\mathcode`l="7160}\displaylimits}%
}

\NoChangel{\log}
\NoChangel{\ln}
\NoChangel{\lim}
\NoChangel{\limsup}
\NoChangel{\liminf}
\NoChangel{\varlimsup}
\NoChangel{\varliminf}
\NoChangel{\varinjlim}
\NoChangel{\varprojlim}

\newenvironment{enumeratea}
{\bgroup\def\theenumi{\alph{enumi}}
\begin{enumerate}}
{\end{enumerate}\egroup}

\newenvironment{smallpmatrix}{\big(\begin{smallmatrix}}{\end{smallmatrix}\big)}
\newenvironment{Smallpmatrix}{\Big(\begin{smallmatrix}}{\end{smallmatrix}\Big)}

\let\ts\textstyle
\let\epsilon\varepsilon

\newcommand{\psfrac}[2]{\sfrac{(#1)}{#2}}
\newcommand{\spfrac}[2]{\sfrac{#1}{(#2)}}
\newcommand{\Psfrac}[2]{(\sfrac{#1}{#2})}

\newcommand\mto{\mathchoice{\longmapsto}{\mapsto}{\mapsto}{\mapsto}}

\usepackage{thmtools, thm-restate}

\newtheorem{theoremIntro}{Theorem}
\newtheorem{theorem}{Theorem}[section]
\newtheorem*{theorem*}{Theorem}
\newtheorem{corollary}[theorem]{Corollary}
\newtheorem{corollaryIntro}[theoremIntro]{Corollary}
\newtheorem{proposition}[theorem]{Proposition}
\newtheorem{propositionIntro}[theoremIntro]{Proposition}
\newtheorem*{conjecture*}{Conjecture}
\newtheorem{lemma}[theorem]{Lemma}

\theoremstyle{definition}
\newtheorem*{definition}{Definition}
\newtheorem*{remark}{Remark}
\newtheorem*{example}{Example}
\newtheorem*{examples}{Examples}

\renewcommand{\H}{\mathcal{H}}
\newcommand{\C}{\mathbb{C}}
\newcommand{\R}{\mathbb{R}}
\newcommand{\N}{\mathbb{N}}
\newcommand{\Z}{\mathbb{Z}}

\newcommand{\HNC}{\mathbb{H}^{n}_{\C}}
\newcommand{\PHNC}{\partial\mathbb{H}^{n}_{\C}}
\newcommand{\HNR}{\mathbb{H}^{n}_{\R}}
\newcommand{\PHNR}{\partial\mathbb{H}^{n}_{\R}}
\newcommand{\PU}{\mathrm{PU}}
\newcommand{\PO}{\mathrm{PO}}
\newcommand{\PNC}{\C \mathbb{P}^{n}}

\DeclareMathOperator{\Id}{Id}
\DeclareMathOperator{\Ricci}{Ricci}
\DeclareMathOperator{\Diag}{Diag}
\newcommand{\tr}{\textup{tr}}


\begin{DefTralics}
\def\C{\mathbb{C}}
\newcommand{\HNC}{\mathbb{H}^{n}_{\C}}
\newcommand{\PU}{\mathrm{PU}}
\end{DefTralics}
